\section{Propagation across a thin cylinder}\label{sec:propagation}

The energy in Proposition~\ref{prop:fixed-barrier} controls the wave on the
circle $|\tau|=2a$, $1\le T\le2$, and is controlled by the data at
$|\tau|=4a$, $T=0$. To compare these two sets, we first carry the energy
out to $|\tau|=1$, where a fixed neighborhood connects the two time
intervals. We then carry it back towards the inner wall. Each passage
costs one logarithm, as Figure~\ref{fig:cylinder} illustrates.
The radial harmonic functions used for comparison are the probabilities
that planar Brownian motion leaves an annulus through one circle or
the other \cite[Theorem~3.18]{MoertersPeres2010}. We use the maximum
principle to retain this comparison in the presence of a potential.

\begin{lemma}\label{lem:cylinder}
Let $0<b<1/4$ and $D=\log(3/b)$. Suppose a continuous potential
$V:[b,3]\to[0,\infty)$ satisfies
\begin{equation}\label{eq:cylinder-moment}
 \int_b^3 s\log(s/b)V(s)\,ds\le M,\qquad
 \sup_{1/2\le s\le3/2} V(s)\le M .
\end{equation}
Let $E\ge0$ be smooth on $\mathbb R_T\times
\{\tau\in\mathbb R^2:b\le|\tau|\le3\}$, with
$\Delta_{T,\tau}E\le V(|\tau|)E$ in the interior. Its circular mean
\[
 H(T,s)=\frac1{2\pi}\int_0^{2\pi}
 E\bigl(T,s(\cos\theta,\sin\theta)\bigr)\,d\theta
\]
satisfies, for $b\le r\le1$,
\begin{equation}\label{eq:cylinder-propagation}
 H(0,r)\ \ge\ c(M)\frac{\log(r/b)}{D^2}\int_1^2 H(T,b)\,dT .
\end{equation}
In particular the factor at $r=2b$ is a constant times $D^{-2}$.
\end{lemma}

\begin{figure}[!htbp]
\centering
\definecolor{cylinderteal}{RGB}{27,88,96}
\definecolor{cylinderorange}{RGB}{153,73,39}
\begin{tikzpicture}[
  x=1cm,y=1cm,font=\small,inner sep=1.5pt,
  wall/.style={draw=black!68,line width=.55pt},
  guide/.style={draw=black!55,line width=.45pt,densely dashed},
  panel/.style={font=\small\bfseries,anchor=west,inner sep=0pt},
  ticklabel/.style={anchor=north,inner ysep=2pt}]
  % The radial coordinates are schematic; the horizontal coordinates give
  % the actual time windows for the two deterministic comparisons.
  \node[panel] at (-.15,2.7) {(a) Outward comparison};
  \fill[cylinderteal!15] (0,1.12) rectangle (4.2,1.58);
  \draw[wall] (0,0) rectangle (4.2,2.05);
  \draw[draw=cylinderteal,densely dashed,line width=.5pt]
    (0,1.35) -- (4.2,1.35);
  \node[fill=cylinderteal!15,inner xsep=3pt,text=cylinderteal]
    at (2.1,1.35) {$3/4\le s\le5/4$};
  \foreach \x/\t in {0/0,1.4/1,2.8/2,4.2/3}{
    \draw[wall] (\x,0) -- +(0,-.09);
    \node[ticklabel] at (\x,-.09) {$\t$};}
  \foreach \y/\s in {0/b,1.35/1,2.05/3}{
    \draw[wall] (0,\y) -- +(-.09,0);
    \node[anchor=east] at (-.11,\y) {$\s$};}
  \node[anchor=south] at (0,2.12) {$s$};
  \node[anchor=west] at (4.32,0) {$T$};
  \draw[draw=cylinderorange,line width=1.8pt] (1.4,0) -- (2.8,0);
  \node[text=cylinderorange,anchor=south] at (2.1,.08) {observed wave};
  \node at (2.1,-.90) {factor at $s=1$: \quad $\displaystyle\frac{\log 3}{\log(3/b)}$};
  \begin{scope}[xshift=6.05cm]
    \node[panel] at (-.15,2.7) {(b) Inward comparison};
    \draw[wall] (0,0) rectangle (4.2,2.05);
    \draw[guide] (0,.57) -- (2.1,.57) -- (2.1,0);
    \foreach \x/\t in {0/{-1},2.1/0,4.2/1}{
      \draw[wall] (\x,0) -- +(0,-.09);
      \node[ticklabel] at (\x,-.09) {$\t$};}
    \foreach \y/\s in {0/b,.57/{2b},2.05/1}{
      \draw[wall] (0,\y) -- +(-.09,0);
      \node[anchor=east] at (-.11,\y) {$\s$};}
    \node[anchor=south] at (0,2.12) {$s$};
    \node[anchor=west] at (4.32,0) {$T$};
    \draw[draw=cylinderteal,line width=1.8pt] (0,2.05) -- (4.2,2.05);
    \node[text=cylinderteal,anchor=north] at (2.1,1.97) {uniform lower bound};
    \fill[cylinderorange] (2.1,.57) circle[radius=2.2pt];
    \node[anchor=west,text=cylinderorange] at (2.23,.57) {$H(0,2b)$};
    \node at (2.1,-.90) {factor at $s=2b$: \quad $\displaystyle\frac{\log 2}{\log(1/b)}$};
  \end{scope}
\end{tikzpicture}
\caption{The two comparisons, with radii $s=|\tau|$ drawn schematically.
The boundary mass on $1\le T\le2$ bounds the mass in the shaded band
from below. The mean inequality on a fixed annulus supplies a lower bound
on the entire upper wall in~(b). The two displayed fractions are values
of the normalized radial harmonic functions. The comparison solutions
with potential dominate fixed multiples of these values, giving the
factor $[\log(3/b)]^{-2}$ at the dot. In the wave estimate, $b=2Kh$:
$H$ controls the wave norm on the marked segment and is controlled by
the initial data at the dot.}\label{fig:cylinder}
\end{figure}


\begin{proof}
Averaging over rotations preserves the differential inequality, so
$H(T,|\tau|)$ is again a supersolution of $(-\Delta_{T,\tau}+V)H\ge0$.
We first estimate the radial solutions that will be used for both
passages. For $S=1$ or $3$ and a fixed $\kappa>0$, consider
\[
 u''+s^{-1}u'=(V+\kappa^2)u,\qquad b<s<S,
\]
with endpoint values $0,1$ in either order. We claim that $u$ is bounded
below by a constant times the function affine in $\log s$ with the same
values. Indeed, $x=\log(s/b)$ changes the equation into
$u_{xx}=P(x)u$ on $[0,A]$, where $A=\log(S/b)$ and
$P(x)=s^2(V(s)+\kappa^2)$. Let $\psi$ solve this equation with
$\psi(0)=0$, $\psi'(0)=1$. Then $\psi(x)\ge x$, and integrating twice gives
\[
 \frac{\psi(x)}x
 =1+\int_0^x\frac{z(x-z)}x P(z)\frac{\psi(z)}z\,dz .
\]
Since $z(x-z)/x\le\min(z,A-z)$, Gronwall's inequality gives
$\psi(A)/A\le\exp(\int_0^A\min(z,A-z)P(z)\,dz)$. The contribution
of $V$ to this integral is at most $M$ by~\eqref{eq:cylinder-moment};
that of $\kappa^2$ is at most
$\kappa^2 S^2\int_0^\infty t e^{-2t}\,dt$, using $t=A-z$.
Thus $\psi(x)/\psi(A)\ge c x/A$. The integral bound is unchanged under
$x\mapsto A-x$, giving the other ordering of the boundary values.

For the outward passage integrate against $\sin(\pi T/3)$ on $[0,3]$.
This cutoff is positive on the observed times and its second derivative
is a multiple of itself. Writing $m(s)=\int_0^3\sin(\pi T/3)H(T,s)\,dT$,
two integrations by parts give
\[
 \int_0^3\sin(\pi T/3)\,\partial_T^2 H(T,s)\,dT
 =-(\pi/3)^2m(s)+(\pi/3)\bigl(H(0,s)+H(3,s)\bigr).
\]
It follows that $m''+s^{-1}m'\le(V+\pi^2/9)m$. Since $m(3)\ge0$,
the maximum principle and the radial estimate with $S=3$ give
$m(s)\ge c\,m(b)\log(3/s)/D$. In particular,
\begin{equation}\label{eq:cylinder-band}
 \int_0^3\int_{3/4\le|\tau|\le5/4}H(T,|\tau|)\,d\tau\,dT
 \ \ge\ c\,m(b)/D .
\end{equation}

On the fixed region $-2<T<4$, $1/2<|\tau|<3/2$, we have
$\Delta H\le MH$. The compact weak Harnack inequality there gives
$H(T,1)\ge c(M)m(b)/D$ for $|T|\le1$. To see the comparison directly,
add a real variable $v$: the function
$H(T,|\tau|)\cos(\sqrt{M+1}\,v)$ is nonnegative and superharmonic on a
fixed slab about $v=0$. The mean inequality on the four-dimensional ball of
radius $2\rho$ at $(x,0)$, integrated over $|v|<\rho$, gives on the
three-dimensional ball of radius $\rho$ about $x=(T,\tau)$ the bound $H(x)\ge c \rho^{-3}\int_{B_\rho(x)}H$ for one small
radius $\rho$ throughout a neighborhood of the compact connected band
$[-1,3]\times\{3/4\le|\tau|\le5/4\}$; here $H(x)=H(T,|\tau|)$.
Cover the band by finitely many balls of radius $\rho/4$. If two such
balls $B,B'$ overlap, then $B'\subset B_\rho(x)$ for every $x\in B$,
so $\inf_B H\ge c \rho^{-3}\int_{B'} H$. A finite chain gives the stated
lower bound from~\eqref{eq:cylinder-band}.

For the return passage take the radial solution with $S=1$,
$\kappa=\pi/2$, $u(b)=0$ and $u(1)=1$. On $|T|<1$, $b<|\tau|<1$,
the function $c(M)m(b)D^{-1}\cos(\pi T/2)u(|\tau|)$ solves
$(-\Delta_{T,\tau}+V)w=0$. It vanishes on $|T|=1$ and $|\tau|=b$,
and the preceding lower bound puts it below $H$ on $|\tau|=1$.
The maximum principle and the radial estimate for $u(r)$ therefore give
\[
 H(0,r)\ge\frac{c(M)m(b)}D\,\frac{\log(r/b)}{\log(1/b)}
 \qquad(b\le r\le1).
\]
This proves~\eqref{eq:cylinder-propagation}, since $\sin(\pi T/3)$ is
bounded below on $[1,2]$.
\end{proof}

\begin{remark}\label{rem:cylinder-power}
The dependence on $r$ and $b$ in Lemma~\ref{lem:cylinder} is sharp
for scalar supersolutions, even when $V=0$. Start with
$(1-\cos T)p(s)$, where $p(s)=\log(3/s)/D$ is the radial harmonic
function with boundary values $1,0$. To cancel the positive part of
its Laplacian without changing those boundary values, put
$g(s)=s^2(1+\log(3/s))/(4D)$. The radial harmonic interpolant
$L(s)=p(s)g(b)+(1-p(s))g(3)$ lies above $g$, since $\Delta_\tau g=p\ge0$.
Thus $E(T,s)=(1-\cos T)p(s)+L(s)-g(s)$ is bounded and nonnegative,
with $\Delta_{T,\tau}E=(\cos T-1)p\le0$. Its trace at $s=b$ is
$1-\cos T$, and since $g$ is increasing we have
\[
 E(0,r)\le (1-p(r))g(3)=\frac94\frac{\log(r/b)}{D^2}.
\]
A fixed positive $r$ would cost one logarithm. The data control $r=2b$,
where a sharper estimate requires more information about the wave energy.
\end{remark}

\begin{proposition}\label{prop:periodic-cap}
On the spatial torus $\TT^n=\RR^n/(100\mathbb Z)^n$, let
$f$ have Fourier support in $c_0/h\le|\xi|\le C_0/h$, $0<h\le1$, and
in one cap of the fixed aperture from Proposition~\ref{prop:fixed-barrier}.
Then
\begin{equation}\label{eq:periodic-cap}
 \left(\int_1^2\|\U(T)f\|_{L^{\pc}(\TT^n)}^{\pc}\,dT\right)^{1/\pc}
 \ \le\ C\,[\log(2/h)]^{2/\pc}\|f\|_{L^{\pc}(\TT^n)} .
\end{equation}
The constant is uniform in $h$ and in the direction of the cap.
\end{proposition}

\begin{proof}
For $f\ne0$, first normalize $\|f\|_{L^2(\TT^n)}=1$.
Fix the constant $K$ in Proposition~\ref{prop:fixed-barrier}, put
$a=Kh$ and $b=2a$, and take $h$ small enough that $b<1/4$. Apply
Lemma~\ref{lem:cylinder} to the energy of that proposition with
$V(s)=Cs^{-2}e^{-cs/h}$. Its first bound in~\eqref{eq:cylinder-moment}
follows from
\[
 \int_b^3 s\log(s/b)V(s)\,ds
 \le C\int_1^\infty \frac{\log t}{t}e^{-2cKt}\,dt ,
\]
and $V$ is uniformly bounded when $1/2\le s\le3/2$. The two energy
bounds in Proposition~\ref{prop:fixed-barrier} are uniform on their
respective circles, so its circular mean satisfies
\[
 H(0,2b)\le C_K a^2\|f\|_{\pc}^{\pc},\qquad
 H(T,b)\ge c_K a^2\|\U(T)f\|_{\pc}^{\pc}.
\]
Take $r=2b$ in~\eqref{eq:cylinder-propagation}, substitute both bounds, cancel
$a^2$, and take $\pc$-th roots. Since $D=\log(3/(2Kh))\lesssim
\log(2/h)$ and $K$ is fixed, this gives~\eqref{eq:periodic-cap}.
Homogeneity removes the normalization after this comparison. For the
remaining $h$ the torus frequencies lie in a fixed ball, where the
estimate follows by a finite Fourier sum.
\end{proof}
